\chapter{Những Điều Về Bạn Bè}
\markboth{Những Điều Về Bạn Bè}{Things a Teacher Wants to Say About Friendship}

\section{Chọn Bạn Tốt}
\begin{truyen}{Chọn Bạn Tốt}{Choose Good Friends}
\chuhoa{B}{ạn bè ảnh hưởng em nhiều hơn em tưởng.} 
\textit{(Friends influence you more than you may realize.)}

Người ở gần mình lâu ngày sẽ kéo nếp sống, cách nghĩ và cả động lực của mình theo.
\textit{(The people around you over time shape your habits, thinking, and motivation.)}

Bạn tốt không nhất thiết là người giỏi nhất lớp. Thường là người tử tế, có trách nhiệm, và không kéo em trượt xuống.
\textit{(A good friend is not necessarily the top student. Often it is someone decent, responsible, and not dragging you downward.)}

Chọn bạn không phải là khinh người khác. Đó là biết giữ mình trước ảnh hưởng lâu dài.
\textit{(Choosing friends is not despising others. It is protecting yourself from lasting influence.)}
\end{truyen}

\begin{baihoc}
Người đi gần em lâu ngày sẽ góp phần tạo nên con người em.
\textit{(Those who walk close to you over time help shape who you become.)}
\end{baihoc}

\begin{ghinhoanh}
\textbf{Short English to remember:}

The people near you help shape your future self.

\end{ghinhoanh}

\begin{tuhocnhanh}
1. Vì sao bạn bè ảnh hưởng mạnh như vậy? \textit{Why do friends influence us so strongly?}

2. Bạn tốt thường có những dấu hiệu nào? \textit{What signs often mark a good friend?}

3. Em đang được kéo lên hay kéo xuống bởi ai? \textit{Who is pulling you up or down right now?}
\end{tuhocnhanh}
\ngancach

\section{Bạn Tốt Giúp Mình Tiến Bộ}
\begin{truyen}{Bạn Tốt Giúp Mình Tiến Bộ}{Good Friends Help Us Improve}
\chuhoa{B}{ạn tốt không chỉ làm mình vui. Họ còn giúp mình tốt lên.}
\textit{(Good friends do not only make us happy. They also help us become better.)}

Có người nhắc em học, có người chỉ em chỗ sai, có người kéo em ra khỏi những ngày chán nản.
\textit{(Some remind you to study, some point out errors, some pull you out of discouraging days.)}

Một tình bạn tốt không phải lúc nào cũng toàn lời dễ nghe. Nó có cả sự thật và sự đỡ đần.
\textit{(A good friendship is not made only of pleasant words. It also contains truth and support.)}

Nếu chơi với người khiến mình đàng hoàng hơn, đó là điều rất đáng quý.
\textit{(If being with someone makes you more decent, that is precious.)}
\end{truyen}

\begin{baihoc}
Bạn tốt không chỉ ở cạnh. Bạn tốt làm em có xu hướng sống và học tử tế hơn.
\textit{(A good friend does not only stay beside you. They incline you toward better living and learning.)}
\end{baihoc}

\begin{ghinhoanh}
\textbf{Short English to remember:}

A good friend helps you become a better self.

\end{ghinhoanh}

\begin{tuhocnhanh}
1. Bạn tốt giúp mình tiến bộ bằng những cách nào? \textit{How do good friends help us improve?}

2. Vì sao tình bạn tốt không chỉ gồm lời dễ nghe? \textit{Why is good friendship not only pleasant words?}

3. Em có đang là kiểu bạn giúp người khác tiến bộ không? \textit{Are you the kind of friend who helps others improve?}
\end{tuhocnhanh}
\ngancach

\section{Đừng Cười Khi Bạn Sai}
\begin{truyen}{Đừng Cười Khi Bạn Sai}{Do Not Laugh When a Friend Is Wrong}
\chuhoa{C}{ười vào cái sai của bạn bè dễ lắm, nhưng hậu quả của nó đôi khi ở lại rất lâu.}
\textit{(Laughing at a friend's mistake is easy, but its effects can last a long time.)}

Một tiếng cười khiến người khác ngượng có thể làm họ ngại nói, ngại học, ngại thử lại.
\textit{(One laugh that embarrasses someone can make them afraid to speak, learn, or try again.)}

Tình bạn tử tế biết dừng lại trước cái sai của nhau bằng sự chừng mực.
\textit{(Decent friendship knows how to meet each other's mistakes with restraint.)}

Bạn mình sai đã đủ khó rồi. Không cần mình làm nó đau thêm.
\textit{(Your friend already feels bad enough when wrong. They do not need you to make it worse.)}
\end{truyen}

\begin{baihoc}
Sự tử tế trong tình bạn thường hiện ra rõ nhất khi người kia đang yếu hoặc đang sai.
\textit{(Kindness in friendship often appears most clearly when the other person is weak or wrong.)}
\end{baihoc}

\begin{ghinhoanh}
\textbf{Short English to remember:}

A kind friend does not add shame to another person's mistake.

\end{ghinhoanh}

\begin{tuhocnhanh}
1. Vì sao cười khi bạn sai dễ làm họ im lại? \textit{Why can laughing at a friend's mistake silence them?}

2. Tình bạn tử tế phản ứng với cái sai ra sao? \textit{How does kind friendship respond to mistakes?}

3. Em đã từng làm ai xấu hổ vì sai chưa? \textit{Have you ever made someone feel ashamed for being wrong?}
\end{tuhocnhanh}
\ngancach

\section{Giúp Bạn Khi Họ Cần}
\begin{truyen}{Giúp Bạn Khi Họ Cần}{Help Friends When They Need It}
\chuhoa{K}{hông phải lúc nào mình cũng giải quyết được vấn đề của bạn. Nhưng có mặt đúng lúc đã là một sự giúp rất thật.}
\textit{(We cannot solve every problem for a friend. But being present at the right time is already real help.)}

Có lúc bạn cần người chỉ bài. Có lúc họ chỉ cần một người không bỏ đi khi họ đang rối.
\textit{(Sometimes a friend needs academic help. Sometimes they just need someone who does not leave while they are struggling.)}

Giúp không nhất thiết phải lớn.
\textit{(Help does not need to be dramatic.)}

Nhiều khi chỉ là nghe kỹ, nhắn hỏi, hay ở cạnh tử tế một chút.
\textit{(Often it is simply listening carefully, checking in, or staying beside them kindly.)}
\end{truyen}

\begin{baihoc}
Có mặt tử tế trong lúc người khác cần là một dạng giúp đỡ rất quý mà nhiều người quên mất.
\textit{(Kind presence in another person's hard moment is a precious form of help that many forget.)}
\end{baihoc}

\begin{ghinhoanh}
\textbf{Short English to remember:}

Sometimes presence is part of help.

\end{ghinhoanh}

\begin{tuhocnhanh}
1. Vì sao có mặt đúng lúc cũng là giúp? \textit{Why is timely presence also help?}

2. Khi bạn cần, em có thể giúp theo những cách nhỏ nào? \textit{How can you help a friend in small ways?}

3. Em có đang bỏ sót ai cần mình hỏi thăm không? \textit{Is there someone you are forgetting to check on?}
\end{tuhocnhanh}
\ngancach

\section{Không Nói Xấu Sau Lưng}
\begin{truyen}{Không Nói Xấu Sau Lưng}{Do Not Speak Badly Behind Someone's Back}
\chuhoa{N}{ói xấu sau lưng thường bắt đầu rất nhẹ: vài câu chê, vài chuyện kể thêm, vài lời nghe cho vui.}
\textit{(Backbiting often begins lightly: a few complaints, a few extra details, a few words for amusement.)}

Nhưng nó làm bầu không khí giữa bạn bè bẩn dần.
\textit{(But it slowly dirties the atmosphere among friends.)}

Người nói xấu người khác dễ làm mất lòng tin của cả tập thể, vì ai rồi cũng sợ mình sẽ là người tiếp theo.
\textit{(Someone who speaks badly about others easily erodes trust, because everyone fears becoming the next target.)}

Nếu có điều cần góp ý, nói thẳng và tử tế vẫn tốt hơn nói sau lưng.
\textit{(If something truly needs to be addressed, speaking directly and kindly is better than talking behind someone's back.)}
\end{truyen}

\begin{baihoc}
Một tập thể khó tử tế nếu mọi người quen sống bằng chuyện nói sau lưng nhau.
\textit{(A group cannot stay decent if people become used to speaking behind one another's backs.)}
\end{baihoc}

\begin{ghinhoanh}
\textbf{Short English to remember:}

Backbiting slowly destroys trust.

\end{ghinhoanh}

\begin{tuhocnhanh}
1. Vì sao nói xấu sau lưng làm mất lòng tin? \textit{Why does backbiting damage trust?}

2. Góp ý thẳng khác với nói sau lưng ra sao? \textit{How is direct feedback different from gossip?}

3. Em có đang giữ một thói quen nói cho vui nhưng làm hại người khác không? \textit{Are you keeping a harmful habit of speaking for amusement?}
\end{tuhocnhanh}
\ngancach

\section{Tôn Trọng Sự Khác Biệt}
\begin{truyen}{Tôn Trọng Sự Khác Biệt}{Respect Differences}
\chuhoa{K}{hông phải bạn nào cũng giống em trong tính cách, tốc độ học, sở thích hay hoàn cảnh sống.}
\textit{(Not every friend is like you in personality, learning speed, interests, or background.)}

Nếu chỉ quý người giống mình, em sẽ rất khó lớn lên trong một thế giới nhiều khác biệt.
\textit{(If you only value people who are like you, it becomes hard to grow in a diverse world.)}

Tôn trọng sự khác biệt không có nghĩa là phải thích mọi thứ. Nó là không coi thường người khác chỉ vì họ khác mình.
\textit{(Respecting differences does not mean liking everything. It means not looking down on others just because they differ from you.)}

Từ sự tôn trọng đó, tình bạn sẽ rộng và lành hơn nhiều.
\textit{(From that respect, friendship becomes broader and healthier.)}
\end{truyen}

\begin{baihoc}
Người biết tôn trọng khác biệt thường sống nhẹ hơn với người khác và cũng bớt cứng với chính mình.
\textit{(People who respect differences often live more gently with others and less rigidly with themselves.)}
\end{baihoc}

\begin{ghinhoanh}
\textbf{Short English to remember:}

Difference is not a reason for contempt.

\end{ghinhoanh}

\begin{tuhocnhanh}
1. Vì sao cần học tôn trọng sự khác biệt từ sớm? \textit{Why should we learn to respect differences early?}

2. Tôn trọng khác với thích hết mọi thứ ở đâu? \textit{How is respect different from liking everything?}

3. Em có đang coi thường ai chỉ vì họ khác mình không? \textit{Are you looking down on someone just because they are different?}
\end{tuhocnhanh}
\ngancach

\section{Học Từ Bạn Bè}
\begin{truyen}{Học Từ Bạn Bè}{Learn from Friends}
\chuhoa{B}{ạn bè không chỉ là người chơi cùng. Họ còn có thể là những người dạy mình rất nhiều điều.}
\textit{(Friends are not only people to spend time with. They can also teach us many things.)}

Có bạn giỏi cách học, có bạn giỏi bình tĩnh, có bạn giỏi sống tử tế trong những chuyện rất nhỏ.
\textit{(Some are good at studying, some at staying calm, some at small acts of kindness.)}

Nếu em biết nhìn, bạn bè quanh em là một lớp học rất sống.
\textit{(If you know how to look, your friends are a living classroom around you.)}

Học từ bạn không làm em kém đi. Nó làm em rộng ra.
\textit{(Learning from friends does not make you smaller. It makes you wider.)}
\end{truyen}

\begin{baihoc}
Biết học từ người gần mình là một dấu hiệu của đầu óc không kiêu và còn muốn lớn.
\textit{(Knowing how to learn from those near you is a sign of a humble mind still willing to grow.)}
\end{baihoc}

\begin{ghinhoanh}
\textbf{Short English to remember:}

Friends can become teachers in everyday life.

\end{ghinhoanh}

\begin{tuhocnhanh}
1. Bạn bè có thể dạy mình những gì? \textit{What can friends teach us?}

2. Vì sao học từ bạn không làm mình nhỏ đi? \textit{Why does learning from friends not reduce us?}

3. Em đã học được điều tốt nào từ một người bạn? \textit{What good thing have you learned from a friend?}
\end{tuhocnhanh}
\ngancach

\section{Cạnh Tranh Lành Mạnh}
\begin{truyen}{Cạnh Tranh Lành Mạnh}{Compete in a Healthy Way}
\chuhoa{C}{ạnh tranh không xấu nếu nó làm em muốn cố hơn chứ không làm em xấu đi.}
\textit{(Competition is not bad if it pushes you to improve rather than corrupting you.)}

Cạnh tranh lành mạnh khiến em nhìn người giỏi hơn như một động lực để học, chứ không như cái gai để ghét.
\textit{(Healthy competition helps you see stronger people as motivation to learn, not as targets of resentment.)}

Nó vẫn giữ được sự tôn trọng với người khác và sự tử tế với chính mình.
\textit{(It preserves respect for others and decency toward yourself.)}

Khi cạnh tranh làm em ghen ghét, chơi xấu hoặc tự hạ mình quá mức, nó đã không còn lành nữa.
\textit{(When competition turns into jealousy, sabotage, or self-contempt, it is no longer healthy.)}
\end{truyen}

\begin{baihoc}
Cạnh tranh tốt làm em cố hơn mà không cần làm ai thấp đi.
\textit{(Good competition makes you try harder without needing to push anyone down.)}
\end{baihoc}

\begin{ghinhoanh}
\textbf{Short English to remember:}

Healthy competition sharpens effort without damaging character.

\end{ghinhoanh}

\begin{tuhocnhanh}
1. Khi nào cạnh tranh còn lành mạnh? \textit{When is competition still healthy?}

2. Dấu hiệu nào cho thấy cạnh tranh đã méo? \textit{What shows that competition has become unhealthy?}

3. Em đang cạnh tranh theo kiểu nào? \textit{What kind of competition are you living in?}
\end{tuhocnhanh}
\ngancach

\section{Đừng Ganh Tị}
\begin{truyen}{Đừng Ganh Tị}{Do Not Let Envy Rule You}
\chuhoa{G}{anh tị làm con người vừa mệt vừa hẹp.}
\textit{(Envy makes people both tired and small.)}

Khi ganh tị, em không còn nhìn người khác bằng con mắt học hỏi nữa.
\textit{(When envy rules, you stop looking at others as people to learn from.)}

Em chỉ thấy điều mình thiếu và điều người khác đang có.
\textit{(You see only what you lack and what others possess.)}

Nếu không cẩn thận, ganh tị sẽ lấy mất cả niềm vui lẫn sự tử tế của em.
\textit{(If you are not careful, envy can take away both joy and decency.)}
\end{truyen}

\begin{baihoc}
Thay vì ganh tị, hãy hỏi mình có thể học gì từ điều tốt của người khác.
\textit{(Instead of envying, ask what you can learn from what is good in others.)}
\end{baihoc}

\begin{ghinhoanh}
\textbf{Short English to remember:}

Envy narrows the heart; learning opens it again.

\end{ghinhoanh}

\begin{tuhocnhanh}
1. Vì sao ganh tị làm người ta mệt? \textit{Why does envy exhaust people?}

2. Em có thể đổi ganh tị thành học hỏi bằng cách nào? \textit{How can envy be turned into learning?}

3. Điều gì ở người khác đang khiến em ganh? \textit{What in others is stirring envy in you?}
\end{tuhocnhanh}
\ngancach

\section{Bạn Bè Là Một Phần Của Trường Học}
\begin{truyen}{Bạn Bè Là Một Phần Của Trường Học}{Friends Are Part of School Life}
\chuhoa{T}{rường học không chỉ gồm bài vở, điểm số và kiểm tra. Nó còn gồm những người đi qua tuổi trẻ cùng em.}
\textit{(School is not only lessons, grades, and tests. It also includes the people who pass through youth with you.)}

Bạn bè làm cho ký ức trường học có hơi ấm.
\textit{(Friends give warmth to school memories.)}

Họ có thể làm em mạnh lên, vui hơn, hiểu đời hơn, hoặc đôi khi làm em đau và học được một bài khác.
\textit{(They can make you stronger, happier, wiser, or at times hurt you and teach a different lesson.)}

Dù thế nào, tình bạn vẫn là một phần rất thật của việc đi học.
\textit{(Either way, friendship remains a very real part of school life.)}
\end{truyen}

\begin{baihoc}
Đi học không chỉ để lấy kiến thức. Em còn đang học cách sống với người khác trong chính những tình bạn của mình.
\textit{(We do not go to school only for knowledge. We are also learning how to live with others through friendship itself.)}
\end{baihoc}

\begin{ghinhoanh}
\textbf{Short English to remember:}

Friendship is also part of education.

\end{ghinhoanh}

\begin{tuhocnhanh}
1. Vì sao bạn bè là một phần của trường học chứ không chỉ là phần bên lề? \textit{Why are friendships part of school itself, not just the side story?}

2. Tình bạn ở trường đang dạy em điều gì? \textit{What is school friendship teaching you?}

3. Em muốn trở thành kiểu bạn nào trong quãng đời học sinh? \textit{What kind of friend do you want to become in school?}
\end{tuhocnhanh}
\ngancach
